%---------------------------------------------------------------------
% The sixteen selected questions.
%
% Every question is reproduced from the end-of-unit Self-Assessment
% Questions of the official AIOU book (1431/5403, 1st ed. 2013); the
% section number after each is where it appears there.
%---------------------------------------------------------------------

\section*{یونٹ \textenglish{1}: \textenglish{ICT} کا تعارف}
\markboth{یونٹ 1: ICT کا تعارف}{}

\begin{sawal}{سوال \textenglish{1} \ \textnormal{\small(\textenglish{1.8, Q.1})}}
\textenglish{Information and Communication Technology} کے فوائد اور نقصانات
تحریر کیجیے۔
\end{sawal}

\begin{hal}
\textbf{تعریف۔} \textenglish{ICT} ان تمام ٹیکنالوجیوں کا مجموعہ ہے جو معلومات
کو جمع کرنے، محفوظ کرنے، عمل میں لانے اور ایک جگہ سے دوسری جگہ منتقل کرنے کے
لیے استعمال ہوتی ہیں --- یعنی کمپیوٹر، انٹرنیٹ، موبائل فون، ریڈیو، ٹی وی اور
نیٹ ورک۔

\medskip
\begin{center}
\begin{tabular}{@{}p{7.0cm} p{7.0cm}@{}}
\toprule
\textbf{فوائد} & \textbf{نقصانات} \\
\midrule
رابطہ تیز اور سستا --- \textenglish{email}، ویڈیو کال & روزگار کا خاتمہ: خودکار نظام سے ملازمتیں کم ہوتی ہیں \\
معلومات تک فوری رسائی، وقت اور سفر کی بچت & نجی معلومات کی عدم تحفظ (\textenglish{privacy}) اور \textenglish{hacking} \\
فاصلاتی تعلیم اور \textenglish{e-learning} ممکن & انحصار: نظام بند ہو تو کام رک جاتا ہے \\
کاروبار میں \textenglish{e-commerce} اور آن لائن بینکاری & صحت پر اثر --- نظر، کمر اور بیٹھے رہنے کے امراض \\
دفتری کام میں درستی اور رفتار & جعلی خبریں اور غیر اخلاقی مواد کا پھیلاؤ \\
دستاویزات کی محفوظ اور کم جگہ گھیرنے والی نقول & ابتدائی لاگت اور تربیت کا خرچ \\
\bottomrule
\end{tabular}
\end{center}

\jawab{\textenglish{ICT} معلومات کو جمع، محفوظ، پروسیس اور منتقل کرنے کی
ٹیکنالوجی ہے۔ بڑے فوائد: تیز رابطہ، معلومات تک رسائی، فاصلاتی تعلیم،
\textenglish{e-commerce}۔ بڑے نقصانات: بے روزگاری، \textenglish{privacy} کا
خطرہ، انحصار، صحت پر اثر۔}
\end{hal}

\begin{sawal}{سوال \textenglish{2} \ \textnormal{\small(\textenglish{1.8, Q.2})}}
تعلیم و تدریس (\textenglish{teaching and learning}) میں \textenglish{ICT} کا
کردار بیان کیجیے۔
\end{sawal}

\begin{hal}
\textbf{استاد کے لیے:}
\begin{itemize}\setlength{\itemsep}{1pt}
  \item \textenglish{multimedia} اور \textenglish{presentation} سے مشکل تصور
        آسانی سے سمجھایا جا سکتا ہے۔
  \item انٹرنیٹ سے تازہ ترین مواد اور مثالیں دستیاب ہوتی ہیں۔
  \item امتحان، نتائج اور حاضری کا ریکارڈ کمپیوٹر پر رکھنا آسان ہے۔
\end{itemize}

\textbf{طالب علم کے لیے:}
\begin{itemize}\setlength{\itemsep}{1pt}
  \item \textenglish{e-learning} اور آن لائن لیکچر سے گھر بیٹھے تعلیم ---
        اے آئی او یو کا پورا نظام اسی پر قائم ہے۔
  \item \textenglish{e-library} اور تحقیقی مواد تک رسائی۔
  \item اپنی رفتار سے سیکھنے کی سہولت؛ لیکچر دوبارہ دیکھا جا سکتا ہے۔
  \item \textenglish{simulation} کے ذریعے تجربہ گاہ کے تجربات کی نقل۔
\end{itemize}

\textbf{ادارے کے لیے:} داخلہ، فیس، نتائج اور رابطے کا خودکار نظام؛ دور دراز
علاقوں تک تعلیم کی رسائی۔

\jawab{\textenglish{ICT} تدریس میں \textenglish{multimedia} اور انٹرنیٹ کے
ذریعے وضاحت آسان بناتی ہے، سیکھنے میں \textenglish{e-learning}، آن لائن لیکچر
اور \textenglish{e-library} فراہم کرتی ہے، اور ادارے کے انتظامی امور کو خودکار
کرتی ہے۔ فاصلاتی تعلیم کا پورا تصور \textenglish{ICT} ہی پر کھڑا ہے۔}
\end{hal}

%---------------------------------------------------------------------
\section*{یونٹ \textenglish{2}: کمپیوٹر کی ساخت}
\markboth{یونٹ 2: کمپیوٹر کی ساخت}{}

\begin{sawal}{سوال \textenglish{3} \ \textnormal{\small(\textenglish{2.9.5, Q.1})}}
کمپیوٹر کیا ہے؟ کمپیوٹر کے پانچ بنیادی اعمال بیان کیجیے۔
\end{sawal}

\begin{hal}
\textbf{تعریف۔} کمپیوٹر ایک الیکٹرانک آلہ ہے جو صارف سے \textbf{ڈیٹا} وصول
کرتا ہے، ہدایات (\textenglish{program}) کے مطابق اس پر عمل کرتا ہے، اور نتیجہ
بطور \textbf{معلومات} پیش کرتا ہے۔

\medskip
\textbf{پانچ بنیادی اعمال:}
\begin{center}
\begin{tabular}{@{}p{1.0cm} p{3.6cm} p{9.4cm}@{}}
\toprule
 & \textbf{عمل} & \textbf{وضاحت} \\
\midrule
۱ & \textenglish{Input}      & کی بورڈ، ماؤس یا اسکینر سے ڈیٹا اندر لینا \\
۲ & \textenglish{Storage}    & ڈیٹا اور ہدایات کو میموری میں محفوظ کرنا \\
۳ & \textenglish{Processing} & \textenglish{ALU} کے ذریعے حسابی و منطقی عمل \\
۴ & \textenglish{Output}     & نتیجہ مانیٹر یا پرنٹر پر پیش کرنا \\
۵ & \textenglish{Controlling}& \textenglish{CU} کے ذریعے باقی چاروں کی نگرانی و ترتیب \\
\bottomrule
\end{tabular}
\end{center}

\jawab{کمپیوٹر وہ الیکٹرانک آلہ ہے جو ڈیٹا لے کر ہدایات کے مطابق پروسیس کرتا
اور معلومات دیتا ہے۔ پانچ بنیادی اعمال: \textenglish{Input}،
\textenglish{Storage}، \textenglish{Processing}، \textenglish{Output} اور
\textenglish{Controlling}۔}
\end{hal}

\begin{sawal}{سوال \textenglish{4} \ \textnormal{\small(\textenglish{2.9.5, Q.2})}}
\textbf{(الف)} مکمل نام لکھیے: \textenglish{ALU}، \textenglish{CU}،
\textenglish{CPU}۔
\textbf{(ب)} \textenglish{Volatile} اور \textenglish{Non-Volatile} میموری کی
تعریف مثالوں سے کیجیے۔
\end{sawal}

\begin{hal}
\textbf{(الف)}
\begin{center}
\begin{tabular}{@{}p{2.2cm} p{5.6cm} p{6.2cm}@{}}
\toprule
\textbf{مخفف} & \textbf{مکمل نام} & \textbf{کام} \\
\midrule
\textenglish{ALU} & \textenglish{Arithmetic and Logic Unit} & جمع، تفریق اور منطقی موازنہ \\
\textenglish{CU}  & \textenglish{Control Unit} & ہدایات کی ترتیب اور دیگر حصوں کی نگرانی \\
\textenglish{CPU} & \textenglish{Central Processing Unit} & \textenglish{ALU + CU +} رجسٹر --- کمپیوٹر کا دماغ \\
\bottomrule
\end{tabular}
\end{center}

\textbf{(ب)}
\begin{center}
\begin{tabular}{@{}p{3.0cm} p{5.6cm} p{5.4cm}@{}}
\toprule
 & \textbf{\textenglish{Volatile}} & \textbf{\textenglish{Non-Volatile}} \\
\midrule
بجلی بند ہونے پر & ڈیٹا ضائع ہو جاتا ہے & ڈیٹا محفوظ رہتا ہے \\
نوعیت & عارضی (\textenglish{temporary}) & مستقل (\textenglish{permanent}) \\
رفتار & زیادہ تیز & نسبتاً سست \\
مثالیں & \textenglish{RAM}، \textenglish{Cache} & \textenglish{ROM}، \textenglish{Hard Disk}، \textenglish{USB}، \textenglish{CD} \\
\bottomrule
\end{tabular}
\end{center}

\jawab{\textenglish{ALU = Arithmetic and Logic Unit}،
\textenglish{CU = Control Unit}، \textenglish{CPU = Central Processing Unit}۔
\textenglish{Volatile} میموری بجلی جاتے ہی خالی ہو جاتی ہے
(\textenglish{RAM})، جبکہ \textenglish{Non-Volatile} میں ڈیٹا محفوظ رہتا ہے
(\textenglish{ROM}، ہارڈ ڈسک)۔}
\end{hal}

\begin{sawal}{سوال \textenglish{5} \ \textnormal{\small(\textenglish{2.9.5, Q.5} اور \textenglish{Q.9})}}
کمپیوٹر کی درجہ بندی مثالوں سے بیان کیجیے، اور پہلی چار نسلوں کا مختصر تعارف
کرائیے۔
\end{sawal}

\begin{hal}
\textbf{(الف) درجہ بندی --- حجم اور صلاحیت کے اعتبار سے:}
\begin{center}
\begin{tabular}{@{}p{3.2cm} p{6.2cm} p{4.8cm}@{}}
\toprule
\textbf{قسم} & \textbf{خصوصیت} & \textbf{استعمال} \\
\midrule
\textenglish{Microcomputer} & ایک صارف، کم قیمت & \textenglish{PC}، لیپ ٹاپ، ٹیبلٹ \\
\textenglish{Minicomputer}  & درمیانی صلاحیت، کئی صارف & چھوٹے ادارے، لیبارٹری \\
\textenglish{Mainframe}     & بہت زیادہ صارف اور ڈیٹا & بینک، ایئر لائن، مردم شماری \\
\textenglish{Supercomputer} & بلند ترین رفتار & موسم کی پیش گوئی، تحقیق، دفاع \\
\bottomrule
\end{tabular}
\end{center}

\textbf{ڈیٹا کی نوعیت کے اعتبار سے:} \textenglish{Analog} (مسلسل مقدار ناپتا
ہے)، \textenglish{Digital} (ہندسوں میں کام کرتا ہے)، اور
\textenglish{Hybrid} (دونوں کا امتزاج --- مثلاً ہسپتال کے مانیٹر)۔

\medskip
\textbf{(ب) پہلی چار نسلیں:}
\begin{center}
\begin{tabular}{@{}p{1.6cm} p{2.6cm} p{4.4cm} p{5.0cm}@{}}
\toprule
\textbf{نسل} & \textbf{دورانیہ} & \textbf{بنیادی پرزہ} & \textbf{خصوصیت / مثال} \\
\midrule
پہلی   & \textenglish{1940--56} & \textenglish{Vacuum Tube}   & بہت بڑا، زیادہ حرارت؛ \textenglish{ENIAC} \\
دوسری  & \textenglish{1956--63} & \textenglish{Transistor}    & چھوٹا اور تیز؛ \textenglish{IBM 1401} \\
تیسری  & \textenglish{1964--71} & \textenglish{IC}            & مزید چھوٹا، کم بجلی؛ \textenglish{IBM 360} \\
چوتھی  & \textenglish{1971--} تا حال & \textenglish{Microprocessor (VLSI)} & ذاتی کمپیوٹر کا آغاز؛ \textenglish{PC} \\
\bottomrule
\end{tabular}
\end{center}

\jawab{درجہ بندی: \textenglish{Micro}، \textenglish{Mini}،
\textenglish{Mainframe}، \textenglish{Super}؛ اور \textenglish{Analog}،
\textenglish{Digital}، \textenglish{Hybrid}۔ نسلیں: \textenglish{Vacuum
Tube}، \textenglish{Transistor}، \textenglish{IC}،
\textenglish{Microprocessor}۔}
\end{hal}

\begin{ghalti}
''تاریخ'' (\textenglish{History}) اور ''نسل'' (\textenglish{Generation}) الگ
چیزیں ہیں۔ تاریخ کمپیوٹر کے ارتقا کا بیان ہے، جبکہ نسل کی بنیاد صرف یہ ہوتی ہے
کہ اُس دور میں کون سا \emph{الیکٹرانک پرزہ} استعمال ہوا۔ جواب میں پرزے کا نام
لکھنا ضروری ہے۔
\end{ghalti}

%---------------------------------------------------------------------
\section*{یونٹ \textenglish{3}: اِن پُٹ آلات}
\markboth{یونٹ 3: ان پٹ آلات}{}

\begin{sawal}{سوال \textenglish{6} \ \textnormal{\small(\textenglish{3.4, Q.1})}}
اسکینر کا مقصد کیا ہے؟ \textenglish{BCR} اور \textenglish{MICR} کو مناسب
مثالوں سے بیان کیجیے۔
\end{sawal}

\begin{hal}
\textbf{اسکینر۔} ایک \textenglish{input device} جو کاغذ پر موجود تصویر یا تحریر
کو پڑھ کر \emph{ڈیجیٹل} شکل میں کمپیوٹر کے اندر منتقل کرتا ہے۔ اس سے دستاویز
دوبارہ ٹائپ کرنے کی ضرورت نہیں رہتی۔

\medskip
\begin{center}
\begin{tabular}{@{}p{2.0cm} p{5.0cm} p{7.4cm}@{}}
\toprule
 & \textbf{مکمل نام} & \textbf{کام اور مثال} \\
\midrule
\textenglish{BCR} & \textenglish{Bar Code Reader}
  & اشیا پر چھپی موٹی پتلی لکیروں (\textenglish{bar code}) کو روشنی کے ذریعے پڑھتا ہے۔ مثال: سپر سٹور پر قیمت اور نام خودکار طور پر آ جانا۔ \\
\textenglish{MICR} & \textenglish{Magnetic Ink Character Recognition}
  & مقناطیسی سیاہی سے لکھے ہندسے پڑھتا ہے۔ مثال: بینک کے چیک کے نیچے چھپا نمبر --- جعل سازی مشکل اور پڑھائی بہت تیز۔ \\
\bottomrule
\end{tabular}
\end{center}

\jawab{اسکینر کاغذی مواد کو ڈیجیٹل شکل میں بدل کر کمپیوٹر میں داخل کرتا ہے۔
\textenglish{BCR} بار کوڈ پڑھتا ہے (سپر سٹور)، اور \textenglish{MICR}
مقناطیسی سیاہی کے حروف پڑھتا ہے (بینک کے چیک)۔}
\end{hal}

\begin{sawal}{سوال \textenglish{7} \ \textnormal{\small(\textenglish{3.4, Q.8})}}
\textenglish{OCR} اور \textenglish{OMR} کا فرق واضح کیجیے۔
\end{sawal}

\begin{hal}
\begin{center}
\begin{tabular}{@{}p{2.6cm} p{5.8cm} p{5.8cm}@{}}
\toprule
 & \textbf{\textenglish{OCR}} & \textbf{\textenglish{OMR}} \\
\midrule
مکمل نام & \textenglish{Optical Character Recognition} & \textenglish{Optical Mark Recognition} \\
کیا پڑھتا ہے & لکھے یا چھپے ہوئے \emph{حروف} & مقررہ خانوں میں لگے \emph{نشان} \\
نتیجہ & قابلِ تدوین متن (\textenglish{editable text}) & نشان کی موجودگی یا غیر موجودگی \\
پیچیدگی & زیادہ --- ہر حرف کی شناخت کرنی پڑتی ہے & کم --- صرف سیاہی کی جگہ دیکھنی ہے \\
مثال & کتاب کا صفحہ اسکین کر کے \textenglish{Word} میں لانا & \textenglish{MCQ} جوابی شیٹ، سروے فارم \\
\bottomrule
\end{tabular}
\end{center}

\jawab{\textenglish{OCR} حروف پہچان کر انہیں قابلِ تدوین متن بناتا ہے، جبکہ
\textenglish{OMR} صرف یہ دیکھتا ہے کہ کسی مقررہ خانے میں نشان لگا ہے یا نہیں۔
\textenglish{OCR} کتاب اسکین کرنے کے لیے، \textenglish{OMR} امتحانی جوابی شیٹ
کے لیے۔}
\end{hal}

%---------------------------------------------------------------------
\section*{یونٹ \textenglish{4}: آؤٹ پُٹ آلات}
\markboth{یونٹ 4: آؤٹ پٹ آلات}{}

\begin{sawal}{سوال \textenglish{8} \ \textnormal{\small(\textenglish{4.4, Q.3})}}
\textenglish{Dot Matrix Printer} اور \textenglish{Inkjet Printer} کا فرق
مثالوں سے واضح کیجیے۔
\end{sawal}

\begin{hal}
پہلے بنیادی تقسیم یاد رکھیے: \textbf{\textenglish{Impact}} پرنٹر کاغذ پر ضرب
لگاتا ہے، \textbf{\textenglish{Non-Impact}} نہیں لگاتا۔

\medskip
\begin{center}
\begin{tabular}{@{}p{2.6cm} p{5.8cm} p{5.8cm}@{}}
\toprule
 & \textbf{\textenglish{Dot Matrix}} & \textbf{\textenglish{Inkjet}} \\
\midrule
قسم & \textenglish{Impact} & \textenglish{Non-Impact} \\
طریقۂ کار & سوئیاں ربن پر ضرب لگا کر نقطے بناتی ہیں & باریک نوزل سے سیاہی کے قطرے پھینکے جاتے ہیں \\
معیار & کم --- نقطے نظر آتے ہیں & بہتر، رنگین چھپائی ممکن \\
آواز & زیادہ & تقریباً خاموش \\
رفتار & سست & نسبتاً تیز \\
فی صفحہ خرچ & بہت کم & زیادہ (\textenglish{cartridge} مہنگا) \\
\textenglish{Carbon copy} & بنا سکتا ہے (ضرب کی وجہ سے) & نہیں بنا سکتا \\
مثال & بینک کی پاس بک، یوٹیلٹی بل & گھر اور دفتر کی عام چھپائی، تصاویر \\
\bottomrule
\end{tabular}
\end{center}

\jawab{\textenglish{Dot Matrix} ایک \textenglish{impact} پرنٹر ہے جو سوئیوں کی
ضرب سے نقطے بناتا ہے --- سستا، شور والا، کم معیار، مگر
\textenglish{carbon copy} بنا سکتا ہے۔ \textenglish{Inkjet}
\textenglish{non-impact} ہے جو سیاہی کے قطرے پھینکتا ہے --- خاموش، رنگین اور
بہتر معیار، مگر فی صفحہ مہنگا۔}
\end{hal}

\begin{sawal}{سوال \textenglish{9} \ \textnormal{\small(\textenglish{4.4, Q.4} اور \textenglish{Q.5})}}
پرنٹر اور پلاٹر میں کیا فرق ہے؟ نیز \textenglish{POS Terminal} اور
\textenglish{ATM} کے افعال بیان کیجیے۔
\end{sawal}

\begin{hal}
\textbf{پرنٹر بمقابلہ پلاٹر:}
\begin{center}
\begin{tabular}{@{}p{2.4cm} p{5.8cm} p{6.0cm}@{}}
\toprule
 & \textbf{\textenglish{Printer}} & \textbf{\textenglish{Plotter}} \\
\midrule
بنیادی کام & متن اور تصویر چھاپنا & لکیروں پر مبنی نقشے اور خاکے بنانا \\
طریقہ & نقطوں (\textenglish{dots}) سے تصویر بنتی ہے & قلم کاغذ پر مسلسل لکیر کھینچتا ہے \\
کاغذ کا حجم & عام طور پر محدود (\textenglish{A4}) & بہت بڑا کاغذ ممکن \\
استعمال & دفتری دستاویزات & نقشہ سازی، \textenglish{engineering drawing}، \textenglish{CAD} \\
\bottomrule
\end{tabular}
\end{center}

\medskip
\textbf{\textenglish{POS Terminal} (\textenglish{Point of Sale})} --- دکان یا
سٹور پر لگا وہ نظام جہاں خریداری مکمل ہوتی ہے۔ یہ \textenglish{bar code}
پڑھتا ہے، قیمت اور کل رقم شمار کرتا ہے، رسید چھاپتا ہے اور
\textenglish{inventory} کا ریکارڈ خودکار طور پر کم کر دیتا ہے۔

\medskip
\textbf{\textenglish{ATM} (\textenglish{Automated Teller Machine})} --- بینک کی
خودکار مشین جو کارڈ اور \textenglish{PIN} کی تصدیق کے بعد رقم نکالنے، رقم جمع
کرانے، بقایا معلوم کرنے اور رقم منتقل کرنے کی سہولت چوبیس گھنٹے دیتی ہے۔

\jawab{پرنٹر نقطوں سے متن و تصویر چھاپتا ہے؛ پلاٹر قلم سے مسلسل لکیریں کھینچ کر
بڑے نقشے بناتا ہے۔ \textenglish{POS} خریداری کے مقام پر قیمت، رسید اور
\textenglish{inventory} سنبھالتا ہے؛ \textenglish{ATM} بینک کے کاؤنٹر کے بغیر
رقم نکالنے، جمع کرنے اور بقایا دیکھنے کی سہولت دیتا ہے۔}
\end{hal}

%---------------------------------------------------------------------
\section*{یونٹ \textenglish{5}: سافٹ ویئر}
\markboth{یونٹ 5: سافٹ ویئر}{}

\begin{sawal}{سوال \textenglish{10} \ \textnormal{\small(\textenglish{5.6, Q.1} اور \textenglish{Q.2})}}
سافٹ ویئر سے کیا مراد ہے؟ \textenglish{System Software} اور
\textenglish{Application Software} میں فرق کیجیے۔
\end{sawal}

\begin{hal}
\textbf{تعریف۔} سافٹ ویئر ہدایات (\textenglish{instructions}) کے اُس مجموعے کا
نام ہے جو کمپیوٹر کو بتاتا ہے کہ کیا کام کرنا ہے۔ ہارڈ ویئر کو چھوا جا سکتا ہے،
سافٹ ویئر کو نہیں۔

\medskip
\begin{center}
\begin{tabular}{@{}p{2.6cm} p{5.8cm} p{5.8cm}@{}}
\toprule
 & \textbf{\textenglish{System Software}} & \textbf{\textenglish{Application Software}} \\
\midrule
مقصد & ہارڈ ویئر کو چلانا اور سنبھالنا & صارف کا مخصوص کام کرنا \\
کس کے لیے & کمپیوٹر خود & صارف (\textenglish{user}) \\
انحصار & اس کے بغیر کمپیوٹر نہیں چلتا & \textenglish{system software} پر چلتا ہے \\
تنصیب & عام طور پر پہلے سے موجود & صارف اپنی ضرورت کے مطابق لگاتا ہے \\
مثالیں & \textenglish{Windows}، \textenglish{Linux}، \textenglish{device drivers}، \textenglish{utilities} & \textenglish{MS Word}، \textenglish{Excel}، براؤزر، گیم \\
\bottomrule
\end{tabular}
\end{center}

\jawab{سافٹ ویئر ہدایات کا مجموعہ ہے۔ \textenglish{System software} ہارڈ ویئر
کو چلاتا اور وسائل سنبھالتا ہے (\textenglish{Windows})، جبکہ
\textenglish{application software} صارف کا مخصوص کام کرتا ہے
(\textenglish{MS Word})۔}
\end{hal}

\begin{sawal}{سوال \textenglish{11} \ \textnormal{\small(\textenglish{5.6, Q.4} اور \textenglish{Q.5})}}
\textenglish{Utility program} کیا ہے اور کیوں استعمال ہوتا ہے؟ نیز
\textenglish{word processing} سے کیا مراد ہے اور اس کی اہم خصوصیات کیا ہیں؟
\end{sawal}

\begin{hal}
\textbf{\textenglish{Utility Program}} --- \textenglish{system software} کی وہ
قسم جو کمپیوٹر کی \emph{دیکھ بھال} کرتی ہے۔ یہ نیا کام پیدا نہیں کرتی بلکہ
نظام کو درست، محفوظ اور تیز رکھتی ہے۔ مثالیں:

\begin{itemize}\setlength{\itemsep}{1pt}
  \item \textenglish{Antivirus} --- وائرس کی نشاندہی اور صفائی
  \item \textenglish{Disk Defragmenter} --- ڈسک پر بکھری فائلیں ترتیب دینا
  \item \textenglish{Backup} اور \textenglish{File Compression}
        (\textenglish{WinZip}) --- نقول اور جگہ کی بچت
  \item \textenglish{Disk Cleanup} --- غیر ضروری فائلوں کا خاتمہ
\end{itemize}

\medskip
\textbf{\textenglish{Word Processing}} --- تحریر کو کمپیوٹر پر لکھنے، درست
کرنے، سجانے، محفوظ کرنے اور چھاپنے کا عمل۔ اہم خصوصیات:

\begin{center}
\begin{tabular}{@{}p{5.2cm} p{9.0cm}@{}}
\toprule
\textenglish{Editing} & متن لکھنا، مٹانا، \textenglish{cut}، \textenglish{copy}، \textenglish{paste} \\
\textenglish{Formatting} & \textenglish{font}، سائز، رنگ، \textenglish{bold}، حاشیے \\
\textenglish{Find and Replace} & کوئی لفظ ڈھونڈ کر پورے مسودے میں بدل دینا \\
\textenglish{Spell \& Grammar Check} & املا اور قواعد کی خودکار جانچ \\
\textenglish{Mail Merge} & ایک خط کو سینکڑوں ناموں کے ساتھ خودکار طور پر تیار کرنا \\
\textenglish{Insert} & تصویر، جدول، چارٹ اور صفحہ نمبر شامل کرنا \\
\bottomrule
\end{tabular}
\end{center}

\jawab{\textenglish{Utility program} نظام کی دیکھ بھال کرنے والا
\textenglish{system software} ہے (\textenglish{antivirus}،
\textenglish{defragmenter}، \textenglish{backup})۔
\textenglish{Word processing} تحریر لکھنے، سجانے اور چھاپنے کا عمل ہے، جس کی
اہم خصوصیات \textenglish{editing}، \textenglish{formatting}،
\textenglish{find \& replace}، \textenglish{spell check} اور
\textenglish{mail merge} ہیں۔}
\end{hal}

%---------------------------------------------------------------------
\section*{یونٹ \textenglish{6}: آپریٹنگ سسٹم}
\markboth{یونٹ 6: آپریٹنگ سسٹم}{}

\begin{sawal}{سوال \textenglish{12} \ \textnormal{\small(\textenglish{6.11, Q.1} اور \textenglish{Q.2})}}
آپریٹنگ سسٹم سے کیا مراد ہے؟ اس کے مختلف افعال بیان کیجیے۔
\end{sawal}

\begin{hal}
\textbf{تعریف۔} آپریٹنگ سسٹم وہ بنیادی \textenglish{system software} ہے جو
کمپیوٹر کے ہارڈ ویئر اور صارف کے درمیان \emph{واسطے} (\textenglish{interface})
کا کام دیتا ہے اور کمپیوٹر کے تمام وسائل کا انتظام کرتا ہے۔ اس کے بغیر کمپیوٹر
استعمال کے قابل نہیں۔ مثالیں: \textenglish{Windows}، \textenglish{Linux}،
\textenglish{macOS}، \textenglish{Android}، \textenglish{UNIX}۔

\medskip
\textbf{اہم افعال:}
\begin{center}
\begin{tabular}{@{}p{4.4cm} p{9.8cm}@{}}
\toprule
\textenglish{Manage Resources} & \textenglish{CPU}، میموری، ڈسک، پرنٹر اور دیگر آلات کی تقسیم و نگرانی \\
\textenglish{User Interface} & صارف سے رابطہ --- آج کل زیادہ تر \textenglish{GUI} کے ذریعے \\
\textenglish{Run Applications} & پروگرام میموری میں لانا، چلانا اور بند کرنا \\
\textenglish{File Management} & فائلیں محفوظ کرنا، نام دینا، \textenglish{folders} کی درخت نما ترتیب \\
\textenglish{Process Management} & کئی پروگراموں کو باری باری \textenglish{CPU} دینا \\
\textenglish{Memory Management} & ہر پروگرام کو میموری دینا اور واپس لینا \\
\textenglish{Security} & \textenglish{password} اور صارف کے اختیارات کی جانچ \\
\textenglish{Built-in Utilities} & \textenglish{antivirus}، \textenglish{backup} جیسے اوزار فراہم کرنا \\
\bottomrule
\end{tabular}
\end{center}

\jawab{آپریٹنگ سسٹم وہ \textenglish{system software} ہے جو صارف اور ہارڈ ویئر
کے درمیان واسطہ بنتا اور تمام وسائل سنبھالتا ہے۔ اہم افعال: وسائل کا انتظام،
\textenglish{user interface}، پروگرام چلانا، فائل، \textenglish{process} اور
میموری کا انتظام، اور سیکیورٹی۔}
\end{hal}

%---------------------------------------------------------------------
\section*{یونٹ \textenglish{7}: ڈیٹا کمیونیکیشن اور نیٹ ورکنگ}
\markboth{یونٹ 7: ڈیٹا کمیونیکیشن اور نیٹ ورکنگ}{}

\begin{sawal}{سوال \textenglish{13} \ \textnormal{\small(\textenglish{7.17, Q.1} اور \textenglish{Q.2})}}
مواصلاتی نظام کے بنیادی اجزا بتائیے، اور \textenglish{Simplex}،
\textenglish{Half Duplex} اور \textenglish{Full Duplex} میں فرق کیجیے۔
\end{sawal}

\begin{hal}
\textbf{مواصلاتی نظام کے پانچ بنیادی اجزا:} پیغام (\textenglish{message})،
بھیجنے والا (\textenglish{sender})، وصول کرنے والا (\textenglish{receiver})،
راستہ (\textenglish{transmission medium}) اور ضابطہ
(\textenglish{protocol})۔

\medskip
\begin{center}
\begin{tabular}{@{}p{2.8cm} p{3.6cm} p{3.4cm} p{4.4cm}@{}}
\toprule
 & \textbf{\textenglish{Simplex}} & \textbf{\textenglish{Half Duplex}} & \textbf{\textenglish{Full Duplex}} \\
\midrule
سمت & صرف ایک طرف & دونوں طرف، مگر باری باری & دونوں طرف، بیک وقت \\
بیک وقت؟ & نہیں & نہیں & ہاں \\
رفتار / کارکردگی & کم & درمیانی & سب سے بہتر \\
مثال & ریڈیو، ٹی وی، کی بورڈ & \textenglish{Walkie-Talkie} & ٹیلی فون، موبائل کال \\
\bottomrule
\end{tabular}
\end{center}

\jawab{اجزا: پیغام، \textenglish{sender}، \textenglish{receiver}،
\textenglish{medium}، \textenglish{protocol}۔ \textenglish{Simplex} میں ڈیٹا
صرف ایک سمت جاتا ہے (ٹی وی)، \textenglish{Half Duplex} میں دونوں سمت مگر باری
باری (\textenglish{walkie-talkie})، اور \textenglish{Full Duplex} میں دونوں
سمت بیک وقت (ٹیلی فون)۔}
\end{hal}

\begin{sawal}{سوال \textenglish{14} \ \textnormal{\small(\textenglish{7.17, Q.5} اور \textenglish{Q.6})}}
\textenglish{LAN} اور \textenglish{WAN} میں فرق کیجیے، اور نیٹ ورک کی مختلف
\textenglish{topologies} بیان کیجیے۔
\end{sawal}

\begin{hal}
\begin{center}
\begin{tabular}{@{}p{2.4cm} p{5.8cm} p{6.0cm}@{}}
\toprule
 & \textbf{\textenglish{LAN}} & \textbf{\textenglish{WAN}} \\
\midrule
مکمل نام & \textenglish{Local Area Network} & \textenglish{Wide Area Network} \\
حدود & ایک عمارت یا کیمپس & شہر، ملک یا پوری دنیا \\
رفتار & زیادہ & نسبتاً کم \\
لاگت & کم & زیادہ \\
ملکیت & عام طور پر ایک ادارہ & کئی ادارے / ٹیلی کام کمپنیاں \\
غلطی کی شرح & کم & زیادہ \\
مثال & دفتر یا لیبارٹری کا نیٹ ورک & انٹرنیٹ، بینک کی ملک گیر شاخیں \\
\bottomrule
\end{tabular}
\end{center}

\medskip
\textbf{\textenglish{Topologies} --- کمپیوٹروں کی ترتیب کا نقشہ:}
\begin{center}
\begin{tabular}{@{}p{2.4cm} p{5.2cm} p{6.6cm}@{}}
\toprule
\textbf{قسم} & \textbf{ترتیب} & \textbf{خوبی / خامی} \\
\midrule
\textenglish{Star} & سب \textenglish{hub} سے جڑے & آسان اور قابلِ اعتماد؛ مگر \textenglish{hub} خراب ہو تو پورا نیٹ ورک بند \\
\textenglish{Ring} & دائرے میں، ہر ایک اگلے سے & ٹکراؤ کم؛ مگر ایک تار ٹوٹنے سے سب متاثر \\
\textenglish{Bus} & سب ایک مرکزی تار سے & سستا اور کم تار؛ مگر مرکزی تار خراب ہو تو سب بند \\
\textenglish{Mesh} & ہر کمپیوٹر ہر دوسرے سے & سب سے زیادہ قابلِ اعتماد؛ مگر بہت مہنگا اور پیچیدہ \\
\bottomrule
\end{tabular}
\end{center}

\jawab{\textenglish{LAN} محدود علاقے کا تیز اور سستا نیٹ ورک ہے،
\textenglish{WAN} وسیع علاقے کا نسبتاً سست اور مہنگا۔ اہم
\textenglish{topologies}: \textenglish{Star}، \textenglish{Ring}،
\textenglish{Bus} اور \textenglish{Mesh}۔}
\end{hal}

\begin{ghalti}
\textenglish{Topology} کے سوال میں صرف نام لکھ دینا کافی نہیں۔ ہر ترتیب کے ساتھ
\textbf{ایک خوبی اور ایک خامی} ضرور لکھیے --- نمبر اسی تقابل کے ملتے ہیں۔ اگر
جگہ ہو تو سادہ خاکہ بھی بنا دیجیے۔
\end{ghalti}

%---------------------------------------------------------------------
\section*{یونٹ \textenglish{8}: ملٹی میڈیا}
\markboth{یونٹ 8: ملٹی میڈیا}{}

\begin{sawal}{سوال \textenglish{15} \ \textnormal{\small(\textenglish{8.9, Q.1} اور \textenglish{Q.4}, \textenglish{Q.5})}}
ملٹی میڈیا سے کیا مراد ہے؟ اس کے اجزا اور اطلاقات تفصیل سے بیان کیجیے۔
\end{sawal}

\begin{hal}
\textbf{تعریف۔} ملٹی میڈیا وہ نظام ہے جس میں معلومات پیش کرنے کے لیے ایک سے
زیادہ ذرائع (\textenglish{media}) کو ملا کر استعمال کیا جاتا ہے --- متن، تصویر،
آواز، ویڈیو اور حرکت۔

\medskip
\textbf{پانچ بنیادی اجزا:}
\begin{center}
\begin{tabular}{@{}p{3.0cm} p{11.2cm}@{}}
\toprule
\textenglish{Text} & بنیادی جزو --- عنوان، وضاحت اور ہدایات \\
\textenglish{Graphics / Images} & تصاویر اور خاکے، جو تصور کو نظر آنے والا بناتے ہیں \\
\textenglish{Audio} & آواز، موسیقی اور تلفظ --- زبان سیکھنے میں خاص طور پر اہم \\
\textenglish{Video} & متحرک حقیقی مناظر \\
\textenglish{Animation} & مصنوعی حرکت --- ایسے عمل دکھانے کے لیے جو حقیقت میں نظر نہ آ سکیں \\
\bottomrule
\end{tabular}
\end{center}

\medskip
\textbf{اہم اطلاقات:} تعلیم و تربیت (\textenglish{CBT}، \textenglish{e-learning})؛
تفریح (فلم، گیم)؛ کاروبار (\textenglish{presentation}، اشتہار)؛
\textenglish{multimedia kiosk} (ہوائی اڈے اور بینک پر معلوماتی سکرین)؛
\textenglish{video conferencing}؛ طب (جسمانی نظام کے متحرک ماڈل)؛ اور
\textenglish{virtual reality}۔

\jawab{ملٹی میڈیا متن، تصویر، آواز، ویڈیو اور \textenglish{animation} کے امتزاج
سے معلومات پیش کرنے کا نظام ہے۔ اجزا یہی پانچ ہیں؛ اطلاقات میں تعلیم، تفریح،
کاروباری \textenglish{presentation}، \textenglish{kiosk} اور
\textenglish{video conferencing} شامل ہیں۔}
\end{hal}

%---------------------------------------------------------------------
\section*{یونٹ \textenglish{9}: کمپیوٹر زبانیں}
\markboth{یونٹ 9: کمپیوٹر زبانیں}{}

\begin{sawal}{سوال \textenglish{16} \ \textnormal{\small(\textenglish{9.15, Q.2} اور \textenglish{Q.5})}}
\textenglish{Low Level} اور \textenglish{High Level Language} میں کیا فرق ہے؟
نیز \textenglish{Compiler} اور \textenglish{Interpreter} میں فرق کیجیے۔
\end{sawal}

\begin{hal}
\begin{center}
\begin{tabular}{@{}p{2.8cm} p{5.6cm} p{6.0cm}@{}}
\toprule
 & \textbf{\textenglish{Low Level}} & \textbf{\textenglish{High Level}} \\
\midrule
قربت & مشین کے قریب & انسان کی زبان کے قریب \\
سمجھنا & مشکل & آسان \\
\textenglish{Portability} & نہیں --- ایک ہی مشین پر چلتی ہے & ہاں --- مختلف مشینوں پر \\
ترجمہ & \textenglish{Assembler} & \textenglish{Compiler} یا \textenglish{Interpreter} \\
رفتار & بہت تیز & نسبتاً کم \\
مثالیں & \textenglish{Machine language}، \textenglish{Assembly} & \textenglish{C}، \textenglish{C++}، \textenglish{Java}، \textenglish{Python} \\
\bottomrule
\end{tabular}
\end{center}

\medskip
\begin{center}
\begin{tabular}{@{}p{2.8cm} p{5.6cm} p{6.0cm}@{}}
\toprule
 & \textbf{\textenglish{Compiler}} & \textbf{\textenglish{Interpreter}} \\
\midrule
ترجمہ & پورا پروگرام ایک ساتھ & ایک وقت میں ایک سطر \\
غلطی کی نشاندہی & ترجمے کے آخر میں، سب اکٹھی & جہاں غلطی آئے وہیں رک جاتا ہے \\
\textenglish{Object code} & بناتا ہے (فائل محفوظ ہو جاتی ہے) & نہیں بناتا \\
دوبارہ چلانا & ترجمہ ایک بار، پھر تیز & ہر بار دوبارہ ترجمہ \\
مجموعی رفتار & تیز & سست \\
مثال & \textenglish{C}، \textenglish{C++} & \textenglish{Python}، \textenglish{BASIC} \\
\bottomrule
\end{tabular}
\end{center}

\jawab{\textenglish{Low level language} مشین کے قریب، تیز مگر مشکل اور غیر
منتقل پذیر ہے؛ \textenglish{high level} انسان کے قریب، آسان اور منتقل پذیر۔
\textenglish{Compiler} پورا پروگرام ایک ساتھ ترجمہ کر کے
\textenglish{object code} بناتا ہے، جبکہ \textenglish{interpreter} ایک ایک
سطر ترجمہ کرتا ہے اور \textenglish{object code} نہیں بناتا۔}
\end{hal}
